% Part: incompleteness
% Chapter: theories-computability
% Section: extensions-of-q-not-decidable

\documentclass[../../../include/open-logic-section]{subfiles}

\begin{document}

\olfileid{inc}{tcp}{cqn}

\olsection{$\Th{Q}$ యొక్క అవైరుధ్య విస్తరణలు అనిర్ణయనీయాలు}

\begin{explain}
ఏ సూత్రం~$!A$కైనా $!A$, $\lnot !A$ రెండింటినీ నిరూపించకపోతే సిద్ధాంతం
\emph{అవైరుధ్యమైనది} అని గుర్తుంచుకోండి. వైరుధ్యం నుంచి ఏదైనా అనుగమిస్తుంది
కనుక, వైరుధ్య సిద్ధాంతం అల్పమైనది: ప్రతి వాక్యమూ నిరూపించదగినదే. ఒక సిద్ధాంతం
$\omega$-అవైరుధ్యమైనదైతే అది అవైరుధ్యమైనదని స్పష్టం. కానీ అవైరుధ్యంగా ఉండడం
బలహీనమైన ఆవశ్యకత (అంటే అవైరుధ్యమైనవే అయినా $\omega$-అవైరుధ్యమైనవి కాని
సిద్ధాంతాలు ఉన్నాయి). మరింత బలమైన సిద్ధాంతాన్ని పొందడానికి
\olref[oqn]{thm:oconsis-q}లోని పరికల్పనను సాధారణ అవైరుధ్యానికి బలహీనపరచవచ్చు.
\end{explain}

\begin{lem}
``సార్వత్రిక గణనీయ సంబంధం'' ఏదీ లేదు. అంటే, కింది ధర్మం గల ద్విస్థానిక గణనీయ
సంబంధం~$R(x,y)$ ఏదీ లేదు: $S(y)$ ఏ ఏకస్థానిక గణనీయ సంబంధమైనా, ప్రతి $y$కు
$S(y)$ సత్యమైతేనే $R(k,y)$ సత్యమయ్యేలా ఏదో~$k$ ఉండాలి.
\end{lem}

\begin{proof}
$R(x,y)$ ఒక సార్వత్రిక గణనీయ సంబంధమని అనుకుందాం. $S(y)$ను $\lnot R(y,y)$
సంబంధంగా తీసుకుందాం. $S(y)$ గణనీయం కనుక ఏదో~$k$కు $S(y)$, $R(k,y)$కు
సమానం. కానీ అప్పుడు $S(k)$ అనేది $R(k,k)$కూ $\lnot R(k,k)$కూ సమానం;
ఇది వైరుధ్యం.
\end{proof}


\begin{thm}
$\Th{Q}$ను కలిగి ఉన్న ఏ అవైరుధ్య సిద్ధాంతమైనా $\Th{T}$గా తీసుకోండి. అప్పుడు
$\Th{T}$ నిర్ణయించదగినది కాదు.
\end{thm}


\begin{proof}
$\Th{T}$ అనేది $\Th{Q}$ యొక్క అవైరుధ్యమైన, నిర్ణయించదగిన విస్తరణ అని
అనుకుందాం. $\Th{T}$ను ఉపయోగించి ఒక సార్వత్రిక గణనీయ సంబంధాన్ని నిర్వచించి
వైరుధ్యాన్ని పొందుతాం.

$R(x,y)$ కిందిది జరిగినప్పుడు, అప్పుడు మాత్రమే వర్తించనివ్వండి:
\begin{quote}
$x$ అనేది $!D(u)$ అనే సూత్రాన్ని సంకేతీకరిస్తుంది; $\Th{T}$ అనేది
$!D(\num y)$ను నిరూపిస్తుంది.
\end{quote}
$\Th{T}$ నిర్ణయించదగినదని ఊహిస్తున్నాం కనుక $R$ గణనీయం. $R$
సార్వత్రికమని చూపుదాం. $S(y)$ ఏ గణనీయ సంబంధమైనా, దానికి $\Th{Q}$లో
(అందువల్ల $\Th{T}$లో కూడా) $!D_S(u)$ అనే సూత్రం ప్రాతినిధ్యం వహిస్తుంది.
అప్పుడు ప్రతి $n$కు
\begin{eqnarray*}
S(\num n) & \lif & T \vdash !D_S(\num n) \\
& \lif & R(\Gn{!D_S(u)}, n)
\end{eqnarray*}
మరియు
\begin{eqnarray*}
\lnot S(\num n) & \lif & T \vdash \lnot !D_S(\num n) \\
& \lif & T \not\vdash !D_S(\num n) \quad \text{($\Th{T}$ అవైరుధ్యమైనది కనుక)} \\
& \lif & \lnot R(\Gn{!D_S(u)}, n).
\end{eqnarray*}
అంటే ప్రతి $y$కు $S(y)$ సత్యమైతేనే $R(\Gn{!D_S(u)}, y)$ సత్యం. కాబట్టి
$R$ సార్వత్రికం; మనం వెతుకుతున్న వైరుధ్యం లభించింది.
\end{proof}

``నిజ అంకగణితం'' అంటే $\Setabs{!A}{ \Sat{\Nat}{!A}}$ అనే సిద్ధాంతం;
అంటే ప్రామాణిక అర్థనిర్దేశంలో సత్యమైన అంకగణిత భాషలోని వాక్యాల సమితి.

\begin{cor}
నిజ అంకగణితం నిర్ణయించదగినది కాదు.
\end{cor}

%In Section~\olref{undefinability:truth:section} we will state a stronger
%result, due to Tarski.

\end{document}
